\section{Tables}

\subsection{Table Between Paragraphs}

Paragraph before the table.

\begin{table}
  \caption{Basic table between two paragraphs.}
  \label{tab:t-basic}
  \begin{autotable}{L R R}
    \toprule
    \textbf{Item} & \textbf{Count} & \textbf{Share} \\
    \midrule
    Alpha & 120 & 0.40 \\ \rowline
    Beta  & 90  & 0.30 \\ \rowline
    Gamma & 90  & 0.30 \\
    \bottomrule
  \end{autotable}
\end{table}

Paragraph after the table.

\subsection{Table Directly After a Heading}

\begin{table}
  \caption{Table right after a heading, with wrapping text and a total row.}
  \label{tab:t-wrap}
  \begin{autotable}{L L R}
    \toprule
    \textbf{Method} & \textbf{Notes} & \textbf{Score} \\
    \midrule
    Baseline & Short note & 0.61 \\ \rowline
    Improved method with a long name & A longer note that wraps onto a second
      line inside the cell & 0.85 \\ \rowline
    \textbf{Total} & & \textbf{1.46} \\
    \bottomrule
  \end{autotable}
\end{table}

\subsection{Automatic Width}

All tables below use the same \texttt{autotable}; only the content differs.

\begin{table}
  \caption{Two short columns: as narrow as their content. This caption is
           long on purpose, to check a long caption above a small table.}
  \begin{autotable}{L R}
    \toprule
    \textbf{Key} & \textbf{Value} \\
    \midrule
    a & 1 \\
    \bottomrule
  \end{autotable}
\end{table}

\begin{table}
  \caption{Medium content: natural width.}
  \begin{autotable}{L R R}
    \toprule
    \textbf{Model} & \textbf{F1} & \textbf{AP} \\
    \midrule
    Model A   & 0.894 & 0.676 \\ \rowline
    Model B   & 0.902 & 0.721 \\
    \bottomrule
  \end{autotable}
\end{table}

\begin{table}
  \caption{Long text in a cell: full width, text wraps.}
  \begin{autotable}{L R}
    \toprule
    \textbf{Observation} & \textbf{Count} \\
    \midrule
    A small object at the same spot in five consecutive samples is missed in every image & 23 \\
    \bottomrule
  \end{autotable}
\end{table}

\subsection{Long Text Columns}

\begin{table}
  \caption{One long text column: it takes most of the width; the short name
           column and the score column stay narrow.}
  \begin{autotable}{L L R}
    \toprule
    \textbf{Method} & \textbf{Notes} & \textbf{Score} \\
    \midrule
    Baseline & Pretrained model used out of the box, keeping only the
      target class; it finds most objects in bright images but almost none in dark ones. & 0.71 \\ \rowline
    Ensemble & Results from both models are merged into one result per object with a
      weighted mean, then linked over time by a tracking step. & 0.91 \\
    \bottomrule
  \end{autotable}
\end{table}

\begin{table}
  \caption{Two text columns of different length: widths follow the amount of
           text in each.}
  \begin{autotable}{L L R}
    \toprule
    \textbf{Idea} & \textbf{Result on val} & \textbf{Kept} \\
    \midrule
    Train on generated labels instead of hand labels & Better on bright images but
      clearly worse on dark ones than training on the hand labels, so the hand
      labels were kept for all later models. & No \\ \rowline
    Half-precision inference & About twice as fast, with slightly lower accuracy. & No \\
    \bottomrule
  \end{autotable}
\end{table}

\subsection{Numbers and Text Together}

\begin{table}
  \caption{Mixed values in one column: set as text (L), left-aligned.}
  \begin{autotable}{L L}
    \toprule
    \textbf{Model} & \textbf{Result} \\
    \midrule
    Model A   & 0.894 (best on dark images) \\ \rowline
    Model B   & 0.902 \\ \rowline
    Model C   & 0.869, but slow (4.4 FPS) \\
    \bottomrule
  \end{autotable}
\end{table}

\begin{table}
  \caption{Better: the number and the note in separate columns, so the numbers
           line up.}
  \begin{autotable}{L R L}
    \toprule
    \textbf{Model} & \textbf{F1} & \textbf{Note} \\
    \midrule
    Model A   & 0.894 & best on dark images \\ \rowline
    Model B   & 0.902 & \\ \rowline
    Model C   & 0.869 & slow (4.4 FPS) \\
    \bottomrule
  \end{autotable}
\end{table}

\subsection{Values Like 0/12/334}

\begin{table}
  \caption{In one column (R): right-aligned, but the parts do not line up.}
  \begin{autotable}{L R}
    \toprule
    \textbf{Model} & \textbf{Found/false/missed} \\
    \midrule
    Baseline  & 0/12/334 \\ \rowline
    Model A   & 194/0/42 \\ \rowline
    Model B   & 15/0/21 \\
    \bottomrule
  \end{autotable}
\end{table}

\begin{table}
  \caption{Better: one column per value under a grouped header, so every part
           lines up.}
  \begin{autotable}{L R R R}
    \toprule
    & \multicolumn{3}{c}{\textbf{Day}} \\
    \cmidrule(lr){2-4}
    \textbf{Model} & \textbf{Found} & \textbf{False} & \textbf{Missed} \\
    \midrule
    Baseline  & 0   & 12 & 334 \\ \rowline
    Model A   & 194 & 0  & 42 \\ \rowline
    Model B   & 15  & 0  & 21 \\
    \bottomrule
  \end{autotable}
\end{table}

\subsection{Narrow Table and Grouped Header}

\begin{table}
  \caption{Grouped header with two groups.}
  \begin{autotable}{L RR RR}
    \toprule
    & \multicolumn{2}{c}{\textbf{Day}} & \multicolumn{2}{c}{\textbf{Night}} \\
    \cmidrule(lr){2-3}\cmidrule(lr){4-5}
    \textbf{Model} & \textbf{Found} & \textbf{Missed} & \textbf{Found} & \textbf{Missed} \\
    \midrule
    A & 194 & 42 & 31 & 5 \\ \rowline
    B & 189 & 47 & 15 & 21 \\
    \bottomrule
  \end{autotable}
\end{table}

\subsection{Many Text Columns}

\begin{table}
  \caption{Four text columns, all with long text: widths are shared out so
           the longest cells wrap to a similar number of lines.}
  \begin{autotable}{L L L L}
    \toprule
    \textbf{Approach} & \textbf{Data used} & \textbf{What happened} & \textbf{Why it was kept or dropped} \\
    \midrule
    Fine-tune the pretrained model on hand labels only &
      All bright and dark samples from the first two weeks, about 1,800 images &
      Dark recall rose from 0.12 to 0.55, bright recall stayed about the same &
      Kept, because it was the first model that worked on dark images at all \\ \rowline
    Add generated labels as extra training data &
      The same samples plus 6,000 images labelled automatically without checking &
      Better on bright images, clearly worse on dark ones, with many false results &
      Dropped, because dark images are the hard case and they got worse \\
    \bottomrule
  \end{autotable}
\end{table}

\begin{table}
  \caption{One long text column and three short ones: the short columns
           stay narrow and the long one takes the rest.}
  \begin{autotable}{L L L L}
    \toprule
    \textbf{Sample} & \textbf{Light} & \textbf{Site} & \textbf{What the model missed and why} \\
    \midrule
    S01 & Dark & Entrance & A dark object lying still in the grass for the whole sequence; it is the same colour as the ground and never moves, so neither the model nor the tracking step picks it up \\ \rowline
    S07 & Bright & Car park & Two objects moving close together are merged into one result when they overlap \\
    \bottomrule
  \end{autotable}
\end{table}

\begin{table}
  \caption{Text and number columns together: text columns wrap, number
           columns keep their natural width.}
  \begin{autotable}{L L L R R}
    \toprule
    \textbf{Model} & \textbf{Training data} & \textbf{Main failure} & \textbf{F1} & \textbf{AP} \\
    \midrule
    Model A & Hand labels from all samples, bright and dark mixed & Small objects far from the camera in dark images & 0.894 & 0.676 \\ \rowline
    Model B & Hand labels plus tracking-filled gaps between labelled images & Two objects that overlap for a long time & 0.902 & 0.721 \\
    \bottomrule
  \end{autotable}
\end{table}

\subsection{Table Longer Than a Page}

A table longer than a page is set where it is written and continues on the
next page, with its header repeated under the caption.

\begin{table}
  \caption{Long table at natural width: it continues on the next page with
           its header repeated.}
  \label{tab:t-long}
  \begin{autotable}{L R R}
    \toprule
    \textbf{Sample} & \textbf{Count} & \textbf{Share} \\
    \midrule
    Sample 01 & 161 & 0.80 \\ \rowline
    Sample 02 & 318 & 0.21 \\ \rowline
    Sample 03 & 229 & 0.82 \\ \rowline
    Sample 04 & 282 & 0.85 \\ \rowline
    Sample 05 & 337 & 0.13 \\ \rowline
    Sample 06 & 350 & 0.06 \\ \rowline
    Sample 07 & 280 & 0.38 \\ \rowline
    Sample 08 & 322 & 0.34 \\ \rowline
    Sample 09 & 138 & 0.65 \\ \rowline
    Sample 10 & 316 & 0.75 \\ \rowline
    Sample 11 & 283 & 0.55 \\ \rowline
    Sample 12 & 367 & 0.24 \\ \rowline
    Sample 13 & 158 & 0.86 \\ \rowline
    Sample 14 & 117 & 0.71 \\ \rowline
    Sample 15 & 239 & 0.06 \\ \rowline
    Sample 16 & 383 & 0.13 \\ \rowline
    Sample 17 & 121 & 0.80 \\ \rowline
    Sample 18 & 61 & 0.43 \\ \rowline
    Sample 19 & 55 & 0.39 \\ \rowline
    Sample 20 & 282 & 0.81 \\ \rowline
    Sample 21 & 238 & 0.59 \\ \rowline
    Sample 22 & 242 & 0.78 \\ \rowline
    Sample 23 & 267 & 0.22 \\ \rowline
    Sample 24 & 227 & 0.17 \\ \rowline
    Sample 25 & 58 & 0.22 \\ \rowline
    Sample 26 & 293 & 0.32 \\ \rowline
    Sample 27 & 172 & 0.91 \\ \rowline
    Sample 28 & 263 & 0.85 \\ \rowline
    Sample 29 & 194 & 0.58 \\ \rowline
    Sample 30 & 299 & 0.54 \\ \rowline
    Sample 31 & 333 & 0.49 \\ \rowline
    Sample 32 & 313 & 0.79 \\ \rowline
    Sample 33 & 248 & 0.79 \\ \rowline
    Sample 34 & 158 & 0.48 \\ \rowline
    Sample 35 & 389 & 0.08 \\ \rowline
    Sample 36 & 183 & 0.82 \\ \rowline
    Sample 37 & 383 & 0.94 \\ \rowline
    Sample 38 & 123 & 0.94 \\ \rowline
    Sample 39 & 207 & 0.74 \\ \rowline
    Sample 40 & 332 & 0.77 \\
    \bottomrule
  \end{autotable}
\end{table}

\Cref{tab:t-long} starts on one page and ends on another.

\subsection{Long Table Directly After a Heading}

\begin{table}
  \caption{Long table with a wrapping text column, directly after a heading.}
  \begin{autotable}{L L R}
    \toprule
    \textbf{Item} & \textbf{Note} & \textbf{Score} \\
    \midrule
    Item 01 & The longest note in the table: it explains what happened in this case and why, so it takes two or three lines. & 0.56 \\ \rowline
    Item 02 & A note of medium length for this row. & 0.95 \\ \rowline
    Item 03 & A longer note that wraps onto a second line inside the cell. & 0.91 \\ \rowline
    Item 04 & Short note. & 0.63 \\ \rowline
    Item 05 & The longest note in the table: it explains what happened in this case and why, so it takes two or three lines. & 0.90 \\ \rowline
    Item 06 & A note of medium length for this row. & 0.86 \\ \rowline
    Item 07 & A longer note that wraps onto a second line inside the cell. & 0.67 \\ \rowline
    Item 08 & Short note. & 0.68 \\ \rowline
    Item 09 & The longest note in the table: it explains what happened in this case and why, so it takes two or three lines. & 0.57 \\ \rowline
    Item 10 & A note of medium length for this row. & 0.54 \\ \rowline
    Item 11 & A longer note that wraps onto a second line inside the cell. & 0.80 \\ \rowline
    Item 12 & Short note. & 0.90 \\ \rowline
    Item 13 & The longest note in the table: it explains what happened in this case and why, so it takes two or three lines. & 0.80 \\ \rowline
    Item 14 & A note of medium length for this row. & 0.55 \\ \rowline
    Item 15 & A longer note that wraps onto a second line inside the cell. & 0.72 \\ \rowline
    Item 16 & Short note. & 0.54 \\ \rowline
    Item 17 & The longest note in the table: it explains what happened in this case and why, so it takes two or three lines. & 0.76 \\ \rowline
    Item 18 & A note of medium length for this row. & 0.59 \\ \rowline
    Item 19 & A longer note that wraps onto a second line inside the cell. & 0.51 \\ \rowline
    Item 20 & Short note. & 0.68 \\ \rowline
    Item 21 & The longest note in the table: it explains what happened in this case and why, so it takes two or three lines. & 0.77 \\ \rowline
    Item 22 & A note of medium length for this row. & 0.99 \\ \rowline
    Item 23 & A longer note that wraps onto a second line inside the cell. & 0.76 \\ \rowline
    Item 24 & Short note. & 0.57 \\
    \bottomrule
  \end{autotable}
\end{table}

\subsection{Table Then Heading}

\begin{table}
  \caption{Table directly followed by a heading.}
  \begin{autotable}{L R}
    \toprule
    \textbf{Key} & \textbf{Value} \\
    \midrule
    x & 1 \\
    \bottomrule
  \end{autotable}
\end{table}
